\documentclass[12pt,a4paper]{report}

% ============================================================
%  GEOMETRY
%  Guide specifies: Left 1.2″, Right 1.0″, Top 1.2″, Bottom 1.0″
% ============================================================
\usepackage[
a4paper,
left=1.2in,
right=1.0in,
top=1.2in,
bottom=1.0in,
headheight=14pt,
headsep=0.2in,
footskip=0.4in
]{geometry}

% ============================================================
%  FONT  – mathptmx gives Times Roman for text AND math
% ============================================================
\usepackage{mathptmx}

% ============================================================
%  LINE SPACING  – guide requires 1.5 everywhere
% ============================================================
\usepackage{setspace}
\onehalfspacing

% ============================================================
%  PARAGRAPH STYLE
% ============================================================
\setlength{\parindent}{0pt}
\setlength{\parskip}{6pt}

% ============================================================
%  ESSENTIAL PACKAGES
% ============================================================
\usepackage{graphicx}
\usepackage{amsmath,amssymb}
\usepackage{booktabs}
\usepackage{multirow}
\usepackage{array}
\usepackage{float}
\usepackage{ragged2e}
\usepackage{caption}
\usepackage{enumitem}
\usepackage{url}
\usepackage{xcolor}
\usepackage[absolute,overlay]{textpos}
\usepackage{tikz}\usepackage{tabularx}
\usepackage{graphicx} % make sure this is in preamble

% ============================================================
%  SVG SUPPORT
%  Compile with: pdflatex --shell-escape
%  Place all .svg files in ./diagrams/
% ============================================================
\usepackage[inkscapelatex=false]{svg}
\IfFileExists{C:/Program Files/Inkscape/bin/inkscape.com}{%
	\svgsetup{inkscapeexe={"C:/Program Files/Inkscape/bin/inkscape.com"}}%
}{%
	\svgsetup{inkscapeexe={inkscape}}%
}
\svgsetup{inkscapepath={svg-inkscape/}}
\svgpath{{./diagrams/}}

% ============================================================
%  FIGURE AND TABLE NUMBERING  – Chapter.Number format (e.g. 3.1, 4.2)
% ============================================================
\renewcommand{\thefigure}{\arabic{chapter}.\arabic{figure}}
\renewcommand{\thetable}{\arabic{chapter}.\arabic{table}}
% Reset counters at each new chapter (default LaTeX behaviour with report class)

% ============================================================
%  HYPERREF  – load LAST among packages
% ============================================================
\usepackage{hyperref}
\hypersetup{colorlinks=false, hidelinks}

% ============================================================
%  TABLE OF CONTENTS / LOT / LOF  (tocloft package)
% ============================================================
\usepackage{tocloft}

% --- TOC title: centred, bold, 16pt ---
\renewcommand{\cfttoctitlefont}{\hfill\bfseries\fontsize{16}{20}\selectfont}
\renewcommand{\cftaftertoctitle}{\hfill}
\setlength{\cftbeforetoctitleskip}{0pt}
\setlength{\cftaftertoctitleskip}{12pt}

% --- Chapter entries in TOC ---
\renewcommand{\cftchappresnum}{CHAPTER\ }
\renewcommand{\cftchapaftersnum}{\quad}
\setlength{\cftchapnumwidth}{7em}
\renewcommand{\cftchapfont}{\bfseries}
\renewcommand{\cftchappagefont}{\bfseries}
\renewcommand{\cftchapleader}{\cftdotfill{\cftdotsep}}
\setlength{\cftbeforechapskip}{6pt}

% --- Section entries ---
\setlength{\cftsecindent}{7em}
\setlength{\cftsecnumwidth}{3em}
\setlength{\cftbeforesecskip}{2pt}

% --- Subsection entries ---
\setlength{\cftsubsecindent}{10em}
\setlength{\cftsubsecnumwidth}{3.5em}
\setlength{\cftbeforesubsecskip}{1pt}

% --- List of Tables title ---
\renewcommand{\cftlottitlefont}{\hfill\bfseries\fontsize{16}{20}\selectfont}
\renewcommand{\cftafterlottitle}{\hfill}

% --- List of Figures title ---
\renewcommand{\cftloftitlefont}{\hfill\bfseries\fontsize{16}{20}\selectfont}
\renewcommand{\cftafterloftitle}{\hfill}

% ============================================================
%  CHAPTER FORMAT
% ============================================================
\usepackage{titlesec}

\renewcommand{\thechapter}{\Roman{chapter}}
\renewcommand{\thesection}{\arabic{chapter}.\arabic{section}}
\renewcommand{\thesubsection}{\arabic{chapter}.\arabic{section}.\arabic{subsection}}

\titleformat{\chapter}[display]
{\centering\bfseries}
{\fontsize{14}{18}\selectfont CHAPTER \thechapter}
{6pt}
{\fontsize{16}{20}\selectfont}
\titlespacing*{\chapter}{0pt}{0pt}{24pt}

% ============================================================
%  SECTION FORMAT
% ============================================================
\titleformat{\section}
{\bfseries\fontsize{14}{18}\selectfont}
{\thesection}{1em}{}
\titlespacing*{\section}{0pt}{12pt}{6pt}

% ============================================================
%  SUBSECTION FORMAT
% ============================================================
\titleformat{\subsection}
{\bfseries\fontsize{12}{15}\selectfont}
{\thesubsection}{1em}{}
\titlespacing*{\subsection}{0pt}{12pt}{6pt}

% ============================================================
%  CAPTION FORMAT
% ============================================================
\DeclareCaptionFont{capfont}{\fontsize{11}{14}\selectfont}
\captionsetup{
	font          = capfont,
	justification = centering,
	skip          = 12pt,
	belowskip     = 6pt
}

% ============================================================
%  REFERENCES ENVIRONMENT
% ============================================================
\newenvironment{myrefs}{%
	\begin{spacing}{1.2}
		\setlength{\parskip}{6pt}
		\setlength{\parindent}{0pt}
		\setlength{\hangindent}{0.31in}
	}{%
	\end{spacing}
}

% ============================================================
%  BIBLIOGRAPHY NAME  – "References"
% ============================================================
\renewcommand{\bibname}{References}

% ============================================================
%  HEADER & FOOTER
% ============================================================
\usepackage{fancyhdr}
\pagestyle{fancy}
\fancyhf{}
\fancyfoot[C]{\thepage}
\renewcommand{\headrulewidth}{0pt}

\fancypagestyle{plain}{
	\fancyhf{}
	\fancyfoot[C]{\thepage}
	\renewcommand{\headrulewidth}{0pt}
}

% To add in toc
\newcommand{\tocentry}[1]{
	\phantomsection
	\addcontentsline{toc}{chapter}{#1}
}


% ============================================================
%  METADATA
% ============================================================
\newcommand{\ReportTitle}{3D Target Shooter Game}
\newcommand{\CourseName}{CSE 4102: Computer Graphics and Image Processing Laboratory}

\newcommand{\FirstAuthorName}{Kazi Rifat Al Muin}
\newcommand{\FirstAuthorRoll}{2107042}

\newcommand{\SubmissionDate}{October, 2026}


\newcommand{\TeacherOne}{Md Tajmilur Rahman}
\newcommand{\TeacherOneTitle}{Lecturer}

\newcommand{\TeacherTwo}{Md. Mubtashim Abrar Nihal}
\newcommand{\TeacherTwoTitle}{Lecturer}

% ============================================================
%  DOCUMENT BEGIN
% ============================================================
\begin{document}

% ----------------------------------------------------------------
%  COVER PAGE
% ----------------------------------------------------------------
\begin{titlepage}
    \pagenumbering{roman}
    \setcounter{page}{1}

    \begin{center}

        {\fontsize{14}{18}\selectfont\textbf{\CourseName}}\\
        \vspace{12pt}

        \begin{spacing}{1.4}
            {\fontsize{18}{27}\selectfont
            \textbf{\textsc{\ReportTitle}}}
        \end{spacing}

        \vspace{2\baselineskip}
        {\fontsize{14}{18}\selectfont By}
        \vspace{2.2\baselineskip}

        {\fontsize{14}{18}\selectfont \textbf{\FirstAuthorName}}\\
        \vspace{5pt}
        {\fontsize{14}{18}\selectfont Roll: \FirstAuthorRoll}

        \vspace{30pt}

        \IfFileExists{images/kuet_logo.png}{
            \includegraphics[height=1.14in,width=1.0in]{images/kuet_logo.png}
        }{
            \fbox{\parbox[c][1.14in][c]{1.0in}{\centering KUET\\Logo}}
        }

        \vspace{1.1\baselineskip}

        {\fontsize{12}{16}\selectfont \textbf{Submitted To}}\\[7pt]

        {\fontsize{12}{15}\selectfont \textbf{\TeacherOne}}\\
        {\fontsize{11}{14}\selectfont \TeacherOneTitle}\\[8pt]

        {\fontsize{12}{15}\selectfont \textbf{\TeacherTwo}}\\
        {\fontsize{11}{14}\selectfont \TeacherTwoTitle}

        \vfill

        \begin{spacing}{1.15}
            {\fontsize{12}{16}\selectfont
            \textbf{Department of Computer Science and Engineering}}\\
            {\fontsize{12}{16}\selectfont
            \textbf{Khulna University of Engineering \& Technology}}\\
            {\fontsize{12}{16}\selectfont
            \textbf{Khulna 9203, Bangladesh}}\\
            {\fontsize{12}{16}\selectfont
            \textbf{\SubmissionDate}}
        \end{spacing}

    \end{center}
\end{titlepage}

\clearpage
\pagenumbering{arabic}
\setcounter{page}{1}

% ============================================================
\section*{Abstract}
% ============================================================
This project develops an end-to-end multilingual sentiment classifier for
KUET-related student social-media posts written in Bangla, English, and
Banglish. The final corpus contains 3,500 posts labeled as
\textit{Positive}, \textit{Negative}, or \textit{Neutral}. The pipeline uses
multilingual text normalization, Banglish-to-Bangla transliteration,
leakage-safe splitting, word and character TF--IDF features, Logistic
Regression, and Transformer fine-tuning. BanglishBERT was evaluated as the
completed Transformer model, while XLM-RoBERTa was configured as an
additional candidate but was not included in the final comparison because
its recorded training run was not completed. Validation Macro-F1 selected
TF--IDF + Logistic Regression, which achieved \textbf{94.29\% test accuracy}
and \textbf{0.9429 Macro-F1} on the untouched 525-post test set.

% ============================================================
\section{Introduction}
% ============================================================
University social-media posts often contain informal spelling, code-mixing,
abbreviations, expressive punctuation, and more than one writing system.
KUET-related posts in this project appear in Bangla, English, and Banglish,
which makes a single monolingual preprocessing strategy unsuitable.
Therefore, the project compares a strong sparse lexical representation with
a pretrained contextual Transformer representation. Model selection is
performed only on a validation set so that the test set remains unbiased.
The resulting system provides three-class sentiment prediction for new
KUET-related posts.

% ============================================================
\section{Objectives}
% ============================================================
\begin{itemize}
    \item To develop a three-class sentiment classifier for KUET-related
    student social-media posts.
    \item To support Bangla, English, and Banglish within one processing
    pipeline.
    \item To preprocess noisy multilingual text while preserving useful
    sentiment information.
    \item To represent text using word/character TF--IDF and contextual
    Transformer embeddings.
    \item To compare TF--IDF + Logistic Regression with BanglishBERT using
    the same train--validation--test split.
    \item To select the final model using validation Macro-F1 and evaluate
    it once on the held-out test set.
    \item To save the trained pipeline and provide a reusable prediction
    function and Streamlit-ready interface.
\end{itemize}

% ============================================================
\section{Dataset and Exploratory Analysis}
% ============================================================
The source Excel sheet contains seven columns:
\texttt{ID}, \texttt{Post Text}, \texttt{Language}, \texttt{Emotion},
\texttt{Sentiment}, \texttt{Source}, and \texttt{Collector}. The classifier
uses \texttt{Post Text} as input and \texttt{Sentiment} as the prediction
target; \texttt{Emotion} and \texttt{Source} are metadata and are not used as
model features. The final classification target has exactly three classes:
Positive, Negative, and Neutral. The dataset contains 3,500 unique posts,
three language categories, an average length of 181.5 characters, and no
exact duplicate post text.

\begin{table}[H]
\centering
\caption{Sentiment distribution by language}
\small
\begin{tabular}{lrrrr}
\toprule
\textbf{Sentiment} & \textbf{Bangla} & \textbf{Banglish} &
\textbf{English} & \textbf{Total} \\
\midrule
Positive & 264 & 284 & 618 & 1166 \\
Negative & 278 & 261 & 628 & 1167 \\
Neutral  & 349 & 371 & 447 & 1167 \\
\midrule
Total & 891 & 916 & 1693 & 3500 \\
\bottomrule
\end{tabular}
\end{table}

The class distribution is almost perfectly balanced. English is the largest
single language group, but Bangla and Banglish together account for 1,807
posts, so multilingual handling is essential.

% ============================================================
\section{Complete Pipeline}
% ============================================================
Figure~\ref{fig:pipeline} summarizes the complete workflow used in the final
project.

\begin{figure}[H]
\centering
\resizebox{0.98\textwidth}{!}{
\begin{tikzpicture}[node distance=0.5cm and 0.85cm]
\node[box] (data) {Excel Corpus\\3,500 Posts};
\node[box, right=of data] (valid) {Validation\\and EDA};
\node[box, right=of valid] (prep) {Normalization +\\Banglish Transliteration};
\node[box, right=of prep] (dedup) {Normalized\\Deduplication};

\node[box, below=0.85cm of valid] (split) {Stratified Split\\70\% / 15\% / 15\%};

\node[model, below left=0.85cm and 0.9cm of split] (tfidf)
{Word + Character\\TF--IDF};
\node[model, below=of tfidf] (lr) {Logistic Regression};

\node[model, below right=0.85cm and 0.9cm of split] (token)
{Subword Tokenization\\+ Dense Embeddings};
\node[model, below=of token] (bert) {BanglishBERT\\Fine-Tuning};

\node[box, below=2.25cm of split] (select)
{Validation Macro-F1\\Model Selection};
\node[box, below=of select] (test)
{Held-Out\\Test Evaluation};
\node[box, below=of test] (deploy)
{Saved Pipeline +\\Prediction Interface};

\draw[arrow] (data) -- (valid);
\draw[arrow] (valid) -- (prep);
\draw[arrow] (prep) -- (dedup);
\draw[arrow] (dedup) |- (split);
\draw[arrow] (split) -- (tfidf);
\draw[arrow] (tfidf) -- (lr);
\draw[arrow] (split) -- (token);
\draw[arrow] (token) -- (bert);
\draw[arrow] (lr) |- (select);
\draw[arrow] (bert) |- (select);
\draw[arrow] (select) -- (test);
\draw[arrow] (test) -- (deploy);
\end{tikzpicture}}
\caption{End-to-end pipeline of the final sentiment-analysis system.}
\label{fig:pipeline}
\end{figure}

% ============================================================
\section{Preprocessing}
% ============================================================
The project uses light preprocessing because aggressive English-only
stemming or stop-word removal can damage Bangla and Banglish information.

\begin{enumerate}
    \item \textbf{Unicode normalization:} every post is converted to NFKC
    form.
    \item \textbf{URL and mention normalization:} links are replaced by
    \texttt{<URL>} and user mentions by \texttt{<USER>}.
    \item \textbf{Noise control:} control characters are removed, extreme
    repetition is capped at three characters, and extra whitespace is
    collapsed.
    \item \textbf{Banglish transliteration:} Banglish rows are converted
    toward Bangla script using Avro phonetic transliteration. Uppercase
    acronyms such as KUET/CSE and recognized English words are protected
    where possible.
    \item \textbf{Two text views:} the Transformer view preserves case,
    while the classical view is lowercased to reduce sparse lexical
    variation.
    \item \textbf{Deduplication:} a normalized duplicate key is produced
    before splitting. No rows were removed in the final corpus because no
    normalized duplicate was found.
\end{enumerate}

The cleaned corpus is divided by stratified sampling with seed 42:
2,450 training posts (70\%), 525 validation posts (15\%), and 525 test posts
(15\%). Each test class contains exactly 175 instances.

% ============================================================
\section{Text Representation and Embedding}
% ============================================================

\subsection{Word and Character TF--IDF}
The final classical branch represents each document using
Term Frequency--Inverse Document Frequency (TF--IDF) \cite{salton}. In
general,

\begin{equation}
\mathrm{TFIDF}(t,d)=\mathrm{TF}(t,d)\,
\log\left(\frac{N}{\mathrm{DF}(t)}\right),
\end{equation}

where $N$ is the number of documents and $\mathrm{DF}(t)$ is the number of
documents containing term $t$. The implementation combines word n-grams
$(1,2)$ with character-within-word n-grams $(3,5)$. Word n-grams capture
phrases, while character n-grams provide robustness to spelling variation
and informal Banglish. The maximum feature counts are 30,000 word features
and 40,000 character features.

\subsection{Contextual Embedding in BanglishBERT}
The Transformer branch does not train Word2Vec on the project corpus.
Instead, BanglishBERT is a pretrained ELECTRA-family discriminator model for
Bengali and English text \cite{banglabert,banglishbert,electra}. Its tokenizer
first divides a post into subword tokens. Each token ID indexes a learned
dense embedding vector,

\begin{equation}
\mathbf{e}_i=E[\mathrm{id}_i],
\end{equation}

and the complete sequence is then contextualized by Transformer
self-attention. For queries $Q$, keys $K$, and values $V$,

\begin{equation}
\mathrm{Attention}(Q,K,V)=
\mathrm{softmax}\left(\frac{QK^T}{\sqrt{d_k}}\right)V.
\end{equation}

Thus, the final representation of a token depends on its surrounding words,
unlike a fixed static word vector. A three-class sequence-classification head
converts the contextual representation into logits for Positive, Negative,
and Neutral. The notebook uses a maximum sequence length of 256 tokens.

% ============================================================
\section{Models and Training}
% ============================================================

\subsection{TF--IDF + Logistic Regression}
The classical model concatenates word and character TF--IDF features and
feeds them to Logistic Regression. For each class $c$,

\begin{equation}
z_c=\mathbf{w}_c^T\mathbf{x}+b_c,\qquad
P(c\mid\mathbf{x})=\frac{e^{z_c}}{\sum_j e^{z_j}}.
\end{equation}

The implementation uses balanced class weighting, a maximum of 3,000
iterations, and seed 42. This model is fast, interpretable at the feature
level, and does not require a GPU.

\subsection{BanglishBERT Fine-Tuning}
BanglishBERT was fine-tuned using a three-class classification head and
class-weighted cross-entropy. The computed training weights were essentially
balanced (approximately 1.001, 1.000, and 1.000), consistent with the nearly
balanced dataset. Training used four epochs, learning rate
$2\times10^{-5}$, physical batch size 4, gradient accumulation of 4
(effective batch size 16), weight decay 0.01, gradient checkpointing, and
FP16 when CUDA was available. The best checkpoint was selected by validation
Macro-F1.

XLM-RoBERTa Base \cite{xlmr} was also listed in the notebook as a second
Transformer candidate. Its recorded run was not completed in the supplied
final artifacts, so no XLM-RoBERTa validation or test score is reported and
it is not used in the final comparison.

% ============================================================
\section{Evaluation Metrics}
% ============================================================
The primary model-selection metric is Macro-F1 because it gives equal
importance to every sentiment class. For class $c$,

\begin{equation}
F1_c =
2\frac{Precision_c\,Recall_c}{Precision_c+Recall_c},
\qquad
MacroF1=\frac{1}{3}\sum_{c=1}^{3}F1_c.
\end{equation}

Accuracy, weighted precision, weighted recall, and weighted F1 are also
reported. The test set is used only after model selection.

% ============================================================
\section{Results}
% ============================================================

The final saved model-comparison results are shown in
Table~\ref{tab:val-results}. The TF--IDF Logistic Regression baseline
achieved the highest validation Macro-F1 and was therefore selected.

\begin{table}[H]
\centering
\caption{Validation model comparison}
\label{tab:val-results}
\small
\begin{tabular}{lcc}
\toprule
\textbf{Model} & \textbf{Accuracy} & \textbf{Macro-F1} \\
\midrule
TF--IDF + Logistic Regression & \textbf{0.9352} & \textbf{0.9353} \\
BanglishBERT & 0.9219 & 0.9219 \\
\bottomrule
\end{tabular}
\end{table}

On the held-out test set, the selected model produced the results in
Table~\ref{tab:test-summary}.

\begin{table}[H]
\centering
\caption{Final held-out test performance}
\label{tab:test-summary}
\small
\begin{tabular}{ccccc}
\toprule
\textbf{Accuracy} & \textbf{Macro-F1} & \textbf{Weighted-F1} &
\textbf{W. Precision} & \textbf{W. Recall} \\
\midrule
0.9429 & 0.9429 & 0.9429 & 0.9431 & 0.9429 \\
\bottomrule
\end{tabular}
\end{table}

\begin{table}[H]
\centering
\caption{Class-wise test performance}
\small
\begin{tabular}{lrrrr}
\toprule
\textbf{Class} & \textbf{Precision} & \textbf{Recall} &
\textbf{F1-Score} & \textbf{Support} \\
\midrule
Positive & 0.9101 & 0.9257 & 0.9178 & 175 \\
Negative & 0.9249 & 0.9143 & 0.9195 & 175 \\
Neutral  & 0.9943 & 0.9886 & 0.9914 & 175 \\
\bottomrule
\end{tabular}
\end{table}

\begin{figure}[H]
\centering
\includegraphics[width=0.66\textwidth]{images/confusion_matrix.jpeg}
\caption{Final test confusion matrix for the selected TF--IDF Logistic
Regression model.}
\label{fig:cm}
\end{figure}

Figure~\ref{fig:cm} shows 162/175 Positive, 160/175 Negative, and 173/175
Neutral posts classified correctly. There were \textbf{30 errors} in total:
15 Negative posts were predicted as Positive, 12 Positive posts as Negative,
one Positive post as Neutral, one Neutral post as Positive, and one Neutral
post as Negative. This is fully consistent with the supplied
\texttt{error\_analysis.csv} and final confusion matrix.

% ============================================================
\section{Inference and Deployment}
% ============================================================
The notebook provides a unified \texttt{predict\_sentiment()} function. For
the selected baseline, a new text is normalized, optionally transliterated
when its language is Banglish, transformed by the saved TF--IDF feature
extractor, and passed to Logistic Regression to obtain three sentiment
probabilities. The trained baseline pipeline, label mapping, metadata,
classification report, model comparison, confusion matrix, and error-analysis
records are saved. The notebook also exports a Streamlit-ready application
where a user can enter a post, select its language, and view the predicted
sentiment and confidence.

% ============================================================
\section{Discussion}
% ============================================================
The results show that the word--character TF--IDF representation is very
effective for this balanced multilingual corpus. Character n-grams help the
model handle Banglish spelling variation, while word n-grams capture short
sentiment-bearing phrases. Banglish transliteration provides a more
consistent script for the Transformer branch, although code-mixed English
words and acronyms must be protected to avoid unwanted conversion.
BanglishBERT achieved a strong validation Macro-F1 of 0.9219, but the
classical baseline achieved the higher validation score of 0.9353.
Because model selection was completed before final test evaluation, choosing
TF--IDF + Logistic Regression did not depend on test performance. Its final
0.9429 Macro-F1 shows strong performance across all three classes, with
Neutral being the easiest class and Positive/Negative confusion accounting
for most remaining errors. The 30 misclassified examples should therefore be
reviewed mainly for mixed sentiment, annotation ambiguity, and
context-dependent expressions.

% ============================================================
\section{Conclusion}
% ============================================================
This project implements a complete multilingual sentiment-analysis pipeline
for KUET student posts in Bangla, English, and Banglish. It combines
multilingual-safe preprocessing, Banglish transliteration, word/character
TF--IDF, Logistic Regression, and Transformer-based comparison. Validation
selected TF--IDF + Logistic Regression as the final model, achieving
94.29\% accuracy and 0.9429 Macro-F1 on the held-out test set. The saved
prediction pipeline and Streamlit-ready interface make the system reusable,
while future work can expand the dataset and complete additional Transformer
experiments such as XLM-RoBERTa.

% ============================================================
%  REFERENCES -- only this bibliography heading is produced
% ============================================================
\begin{thebibliography}{99}
\small

\bibitem{salton}
G. Salton and C. Buckley,
``Term-weighting approaches in automatic text retrieval,''
\textit{Information Processing \& Management},
vol. 24, no. 5, pp. 513--523, 1988,
doi: 10.1016/0306-4573(88)90021-0.

\bibitem{sklearn}
F. Pedregosa \textit{et al.},
``Scikit-learn: Machine learning in Python,''
\textit{Journal of Machine Learning Research},
vol. 12, pp. 2825--2830, 2011.
[Online]. Available:
\url{https://jmlr.org/papers/v12/pedregosa11a.html}

\bibitem{electra}
K. Clark, M.-T. Luong, Q. V. Le, and C. D. Manning,
``ELECTRA: Pre-training text encoders as discriminators rather than
generators,''
in \textit{Proc. Int. Conf. Learning Representations (ICLR)}, 2020.
[Online]. Available:
\url{https://openreview.net/forum?id=r1xMH1BtvB}

\bibitem{banglabert}
A. Bhattacharjee, T. Hasan, W. Ahmad, K. S. Mubasshir, M. S. Islam,
A. Iqbal, M. S. Rahman, and R. Shahriyar,
``BanglaBERT: Language model pretraining and benchmarks for low-resource
language understanding evaluation in Bangla,''
in \textit{Findings of the Association for Computational Linguistics:
NAACL 2022}, pp. 1318--1327, 2022,
doi: 10.18653/v1/2022.findings-naacl.98.

\bibitem{banglishbert}
CSE BUET NLP Group,
``BanglishBERT,'' Hugging Face Model Hub.
[Online]. Available:
\url{https://huggingface.co/csebuetnlp/banglishbert}
[Accessed: Sep. 23, 2026].

\bibitem{xlmr}
A. Conneau \textit{et al.},
``Unsupervised cross-lingual representation learning at scale,''
in \textit{Proc. 58th Annual Meeting of the Association for Computational
Linguistics}, pp. 8440--8451, 2020,
doi: 10.18653/v1/2020.acl-main.747.

\bibitem{transformers}
T. Wolf \textit{et al.},
``Transformers: State-of-the-art natural language processing,''
in \textit{Proc. 2020 Conf. Empirical Methods in Natural Language
Processing: System Demonstrations}, pp. 38--45, 2020.
[Online]. Available:
\url{https://aclanthology.org/2020.emnlp-demos.6/}

\end{thebibliography}

\end{document}
